\preludeandfugue{The delights of dining well}

\contributor{Sade Kaye}

% Key: delectation taste loneliness

\prelude

\paragraph{Dining out}\strut

There was nothing left of the money except the taxi fare home.  Where
had it all gone? Was it worth it?

She'd gone out that evening with her purse stuffed with ten pound
notes.  It was all her recent winnings from her premium bonds and
she'd decided to spend it all in one go on the best meal out she could
ever have.  This would be one of those gourmet taster meals with many
tiny courses, delightful and full of flavour, but each one fleeting and
ephemeral, like the wood smoke from smouldering applewood that
surrounded her dish of smoked salmon, or the aroma of freshly mown
grass trapped under the glass cover of her exquisite but miniscule
steak.

With no-one in her life, she'd gone to dine alone.  This didn't bother
her at all, she told herself several times that evening.  She noted
how her lone pigeon breast in madeira~\textit{jús} was happy on its
own, and so was she. A single but beautifully carved radish sat on top
of her salad of wilted chicory leaves, and it showed her how elegant
it was to be on one's own.  Of course this was not always true, she
realised as she enjoyed a dish consisting of two friendly tortellini
stuffed with wild mushrooms on a creamy Parmesan sauce and topped with
gratings of truffle. But she tasted every dish slowly and carefully,
savouring every tiny mouthful and she beamed beautiful appreciative
smiles to the waiters when they cleared her table and brought the next
dish to her.

This was her special treat to herself. She was not lonely because she
was having the greatest meal of her life.  These dishes being set in
front of her were her friends and would be for ever. She was doing
this for herself because she deserved this extravagant treat, and she
had been lucky enough to win the money from ERNIE. She had appreciated
every moment, she told herself, as she finished the last morsel of her
last dish: a vanilla and cherry tart flavoured with kirsch.

All through the delightful meal she had been counting the dishes as
they arrived and committing the list of courses to memory.  Some day
when he was a bit older, she would tell her nephew Alan all about it.
After all, he was already watching all the food programmes on TV and
had said that he wanted to be a chef when he grew up.  She'd watch his
face intently as she described the flavours, textures and aromas of
each dish, and would watch his reaction with the same pleasure.

She giggled involuntarily as she got up from the table.  `I'll be
dining out on this for the rest of my life,' she said.

\fugue

\paragraph{Take a chance on lunch}\strut

Lunch hasn't treated me right. Yesterday I had soup. The soup was too
salty. The day before I had steak. It was over-cooked and
chewy. Sunday roast was pheasant. It was like the ace in a suit of
cards: high. I'm sure it wasn't meant to be quite that high. What's more, it
was very dry, and I needed more jús, but there was none.

So today I am taking my lunch with trepidation. Here I go again, here
I go, taking a chance...

Today I have ordered lamb chops. The menu promises me young lamb but
knowing my luck I'm sure I'll end up with tired old mutton. I like my
lamb slightly pink. In a possibly unwarranted rash of optimism that's
what I've ordered, but now that I'm sitting at the table waiting, I am
wondering if it was a mistake. After all, undercooked mutton can be
almost indigestible.

It's not that I expect the chef to see every dish in completely
the same way as I do. I am willing to and expect to make compromises. If I
were chef and had to write the menus every day I'm sure I'd run out of
ideas. In any case, I love my taste buds to be challenged with something
beautiful and exciting and new.

Presentation is of course essential too. 
Of course taste matters. But when I want to send a quick photo from my
phone to tell all my friends what a wonderful meal I'm eating, it's no
good saying it tastes good. It must also look good.
Our primary sense is sight,
and if the dish doesn't look good, it can be off-putting. As much as I
like to be politically correct and say the taste is the only thing
that matters, presentation and aroma are very important. However, one
mustn't over-do it. Good looks are necessary but I don't want a
dandified dish that says, `look at me'. Nor one that is over-perfumed
that says, `smell me', but just something attractive and appetising
that says, `eat me'.

And a really good lunch must be just filling enough, as well as fulfilling. I
can't be having a quick snack and then wandering off wondering what to
eat next. I need my lunch to be sufficiently satisfying to
last just long enough
until my next meal. Nor do I want the sort of lunch that requires a 
lie-down afterwards. I know there are those that do, but all that seems a
bit hedonistic to me.  I just want my taste buds and memory entertained all
afternoon until the next meal.

But I must stop wittering. The waiter has just brought over my dish, and I
must give it the best attention I can.

Mmmm\ldots

Yes, it is delicious. Even better than it looks. It smells appetising
and tastes even better. Sweet and tender flesh that melts in the
mouth. The rich red wine j\'us complements it perfectly, and
it is accompanied with three small buttery new potatoes\=steamed
not boiled to preserve all their flavour.  There are also
two delicious and just-crisp spears of asparagus, a saut\'eed
half-shallot and a few leaves of dandelion whose bitterness complements
the sweetness of the j\'us perfectly. This dish is one of those
that belies its simplicity, and every moment of the preparation from the
gentle cooking of the meat, to the sauce and accompanying vegetables
has to be exactly perfect.  It is truly Beautiful.  A work of art.

As much as it has been prepared very skilfully, I know it is a
simple-looking dish, and I am sure some people would make fun of
me for falling for something so plain, but I wasn't looking for
anything ornate and flashy.  I find a fancy dish can only last a short
time on the kitchen menu before one gets tired of it and never wants
it again. But a good dish has infinitely many subtle variations.  One
could change the cut of the lamb, develop the stock, or use red wine
from a a different grape, and of course one should always serve such
artistry only with
vegetables that are the very best of the season.  This is a dish that
will grow and adapt with me for my whole life, and one I'd happily
introduce to my friends (when I find some) for an occasional lunch party.

Yes, I have been lucky in lunch after all. Here is something to
remember and enjoy for the rest of my life.


